\section{A checker is not the same as an explanation}
Lean is a programming language and a \emph{proof assistant}: a program that checks proofs. In Lean, every proposition has a corresponding type, and a proof of the proposition is an object of that type. When you write a proof, Lean checks that the object you built really has the type you claimed. Most proofs are written with \emph{tactics}, which are instructions that build the proof object step by step. The official introductions explain these ideas in more detail~\cite{leanprop,leantactics}.

Our aim is not to memorize a catalogue of tactics. It is to make the discipline of Chapter~4 explicit: which objects and hypotheses are available, and what exactly remains to be established? The examples below reproduce mathematics already learned, so that a difficulty with Lean's notation can be told apart from a difficulty with the mathematics.

To run the examples yourself, install Lean using the official instructions, which use the VS Code editor with the Lean~4 extension~\cite{leaninstall}. The examples in Chapters~15 and~16 use only Lean's built-in libraries, not the large mathematical library Mathlib, and they have been checked with Lean version~4.19.0.

\subsection*{From a mathematical sentence to a file}
A Lean file is an ordinary text file whose name ends in \code{.lean}, for example \code{proofs.lean}. It contains \emph{declarations}: some define objects or meanings, and others state and prove results. All the examples in Chapters~15 and~16 can be collected in one such file. Appendix~\ref{app:practical} shows how to start it. Its first line, \code{import Std}, makes Lean's standard library available. A line \code{namespace MathCourse} then groups our names so that they do not collide with names defined elsewhere, and a closing \code{end MathCourse} ends that group. This \code{end} closes a group of names; it has nothing to do with the end of a proof.

You can check the file in two ways. In the editor, Lean checks the file continuously and shows, beside each line of a proof, what has been established so far and what remains to be proved. From a terminal---a window for typing commands to installed programs---the command \code{lean proofs.lean} asks Lean to check the whole file at once. Lean proof lines are typed into the file, not into the terminal.

You can read the mathematics in this chapter before installing anything. Once Lean is installed, compare the hypotheses and goal that the editor displays after each line with the explanations given here; learning to read those displays is one of the main goals of this chapter.
